\subsubsection{Effet de la divergence à grande longueur}
Le chemin d'édition injecté donne une borne supérieure \(d(A,B)\le e\). Au score \(s\), WFA conserve le plus grand indice de ligne \(F_s(k)\) atteint sur la diagonale \(k=i-j\). Prolonger une correspondance ne consomme aucun point de score~\cite{wfa}.
L'algorithme~\ref{alg:wfa} explicite le cas des coûts unitaires de WFA2-lib.
\begin{algorithm}[!htbp]
\caption{Progression d'un front à coûts unitaires}\label{alg:wfa}
\begin{algorithmic}[1]
\Require Strings \(A[1:m],B[1:n]\)
\Ensure Unit-cost global edit distance \(d(A,B)\)
\State \(s\gets0\); extend the common prefix on diagonal \(k=0\)
\While{the current front does not reach \((m,n)\)}
  \State \(s\gets s+1\)
  \ForAll{valid diagonals \(k\) with \(|k|\le s\)}
    \State Choose the furthest valid insertion, deletion or substitution predecessor
    \State Extend matching letters and store the final row index in \(F_s(k)\)
  \EndFor
\EndWhile
\State \Return \(s\)
\end{algorithmic}
\end{algorithm}
Pour les prédécesseurs intérieurs, les indices de ligne après délétion, insertion ou substitution valent respectivement \(F_{s-1}(k-1)+1\), \(F_{s-1}(k+1)\) et \(F_{s-1}(k)+1\). Chaque mouvement reste dans le rectangle des préfixes. Au plus \(2s+1\) diagonales sont actives au score \(s\), soit \(O((d+1)^2)\) entrées de front jusqu'au score final, auxquelles s'ajoutent les comparaisons des lettres égales.
\chart{comparison/temps_distance.png}{Les durées sont rapportées à la longueur et à la divergence. À gauche, le deuxième appel de chaque processus porte sur la première paire à 5\,\% de chaque longueur. À droite, cinq paires distinctes par taux sont mesurées à froid. Les méthodes CPU tournent sur EPYC et NvPD-CUDA sur GH200.}{fig:temps}{bench/campaign/run.py}
\FloatBarrier
À longueur fixée, NvPD parcourt toujours \(mn\) cases. La figure~\ref{fig:temps} montre que le temps de WFA2 augmente avec la divergence. Le débit dense \(mn/(10^9 T)\) compte des milliards de cases utiles par seconde pour NvPD, hors recalcul des halos. Pour un front, le même quotient est un débit équivalent.
Pour les longueurs 5\,000, 20\,000, 100\,000 et un million, le panneau a) compte respectivement neuf, neuf, cinq et cinq processus par méthode. À un million de bases, le parcours séquentiel prend \(1\,294{,}72\pm83{,}01\), \(1\,291{,}85\pm80{,}46\) et \(1\,264{,}36\pm82{,}90\)\,s à 2, 5 et 10\,\% d'éditions injectées. Les cinq processus de chaque taux travaillent simultanément sur des cœurs distincts, chacun après un alignement d'échauffement. À 5\,\%, l'IC à 95\,\% s'étend de 1\,191,95 à 1\,391,76\,s. La demi-largeur relative atteint 7,7 à 8,1\,\% selon le taux. Les durées se répartissent entre deux groupes de cœurs, sans mesure isolant l'origine matérielle de cet écart.

Sur les cinq paires distinctes de chaque taux à un million de bases, la moyenne et l'écart-type pour WFA2 passent de \(0{,}610\pm0{,}009\) à \(11{,}638\pm0{,}063\)\,s entre 2 et 10\,\% d'éditions injectées. Sur les mêmes groupes, Edlib passe de \(1{,}706\pm0{,}006\) à \(6{,}024\pm0{,}005\)\,s.

Pour comparer les temps sur le même nœud GH200, la figure~\ref{fig:bilan} reprend une paire de chacune des douze configurations. NvPD utilise le GPU, Edlib et WFA2 un cœur Grace. Chaque bloc mesure les trois méthodes sur la même paire après échauffement, hors lecture du FASTA. Les configurations de 5 et 20~knt comptent neuf blocs, celles de 100~knt et 1~Mnt en comptent cinq.
\chart{comparison/synthese.png}{Rapports des temps de NvPD sur GPU à ceux des bibliothèques sur un cœur Grace. Chaque case donne la moyenne des rapports appariés, leur écart-type et l’intervalle à 95\,\% à paire fixée. Un rapport inférieur à un favorise NvPD.}{fig:bilan}{python/studies/score_comparison.py}
À 10\,\% d’éditions introduites, le rapport NvPD/WFA2 passe de 1,054 à 0,365 entre 100~knt et 1~Mnt. À 1~Mnt et 5\,\%, il vaut 1,305. Ces mesures répondent à une question de temps d’attente pour une paire soumise à chacune de ces configurations matérielles. La série suivante reprend chaque configuration sur cinq paires distinctes.

\synthese{À un million de bases, passer de 2 à 10\,\% d’éditions introduites multiplie par dix-neuf le temps moyen de WFA2 sur les cinq paires de chaque taux. Dans la comparaison sur GH200, seule la paire de 1~Mnt à 10\,\% place NvPD devant les deux bibliothèques.}


\subsubsection{Variabilité entre paires à divergence fixée}\label{sec:divergence-paires}
Pour chaque longueur et chaque taux d’éditions introduites, cinq paires synthétiques sont alignées sur le même nœud GH200. Chaque méthode effectue trois blocs par paire. Chaque bloc comprend un échauffement suivi d’un appel chronométré. NvPD utilise le GPU et chaque bibliothèque un cœur Grace. Le chronomètre exclut la lecture des séquences.

Pour une paire \(i\), le rapport \(R_i\) est la moyenne des trois rapports de temps appariés. La figure~\ref{fig:divergence-paires} montre la moyenne et l’écart-type des cinq valeurs \(R_i\). Les intervalles à 95\,\% utilisent la loi de Student à quatre degrés de liberté. Les taux décrivent les opérations introduites dans le générateur. La distance minimale peut être inférieure à leur nombre.
\chart{comparison/divergence_paires.png}{Rapports des temps de NvPD sur GPU à ceux d’Edlib et de WFA2 sur un cœur Grace. Les points donnent la moyenne sur cinq paires. Les traits épais donnent l’écart-type entre paires et les traits fins l’intervalle à 95\,\%. Un rapport inférieur à un favorise NvPD.}{fig:divergence-paires}{bench/repeated/divergence/analyze.py}

À un million de bases et 10\,\% d’éditions introduites, NvPD devance les deux bibliothèques sur les cinq paires. Le rapport NvPD/WFA2 vaut \(0{,}3706\pm0{,}0038\), avec un intervalle de 0,3658 à 0,3754. Le rapport NvPD/Edlib vaut \(0{,}6210\pm0{,}0064\), avec un intervalle de 0,6131 à 0,6289. À 100\,000 bases et 10\,\%, NvPD devance WFA2 sur deux paires. Le rapport moyen vaut \(1{,}044\pm0{,}048\), avec un intervalle de 0,985 à 1,103.

\synthese{À un million de bases, les rapports moyens NvPD/WFA2 valent 4,821 à 2\,\%, 1,309 à 5\,\% et 0,371 à 10\,\%. Le classement dépend de la longueur et du taux d’éditions introduites.}

\subsubsection{Empreinte mémoire selon la longueur et la divergence}
À longueur fixée, WFA2 explore davantage de diagonales lorsque la distance augmente. Dans la version dense initiale de NvPD, la taille des tampons dépend de la longueur des séquences. Le calcul du score place la plus longue chaîne en colonne. Sur l'alphabet à quatre bases, le GPU réserve quatre lignes d'indices et deux lignes de scores, chacune en entiers de quatre octets. Pour des séquences de longueurs \(m,n\), cette allocation vaut
\[
  M_{\mathrm{device}}=4(4+2)\bigl(\max(m,n)+1\bigr)\quad [\mathrm{bytes}].
\]
Le panneau a) de la figure~\ref{fig:synthetic-memory} porte sur les quinze paires de chaque longueur. Le RSS est le pic résident du processus hôte. Les panneaux b) à e) utilisent la première paire de chaque taux. Pour les deux bibliothèques sur Grace, un second processus relève le RSS après lecture et pendant les appels, avec un intervalle cible d'une milliseconde. Le maximum ainsi échantillonné est une borne inférieure du pic de mémoire.

\chart{comparison/memoire_synthetique.pdf}{Le panneau a) compare les tampons CUDA au RSS hôte de NvPD (échelle logarithmique). Il donne la moyenne et les extrêmes sur quinze paires distinctes par longueur. Les panneaux b) à e) donnent les moyennes des maxima échantillonnés pour Edlib et WFA2. Les traits épais indiquent un écart-type, les traits fins l'IC à 95\,\% sur neuf processus à 5 et 20~knt et cinq aux deux autres longueurs. Les échelles horizontales sont propres à chaque longueur.}{fig:synthetic-memory}{python/studies/synthetic_memory.py}
À 1~Mnt, les tampons du noyau initial occupent de 24\,000\,024 à 24\,016\,800 octets sur les quinze paires. Les insertions et les délétions changent leur taille lorsque la chaîne mutée dépasse la référence. Le pic RSS hôte de NvPD vaut 216\,Mio pour chacune de ces paires, contre 180\,Mio de 5 à 100~knt.

Avec le calcul en bande, un groupe de 32 colonnes occupe 32 octets dans l’index. Les quatre masques de présence et les quatre positions précédant le groupe permettent de retrouver la dernière occurrence de chaque base. Les scores occupent deux tranches dont la largeur dépend de la borne \(U\). Le panneau a) de la figure~\ref{fig:nvpd-focus-memoire} compare les allocations explicites des deux noyaux.

Le panneau b) mesure l’effet du nombre de processus. Sur un même GH200, cinq paires de 100\,000 bases à 10\,\% sont traitées successivement par un processus, puis comparées à une exécution concurrente avec un processus par paire. L’ordre des deux modes varie entre les cinq blocs. Après échauffement, Linux attribue à chaque processus une part \(1/k\) de chaque page partagée par \(k\) processus. La somme de ces mémoires résidentes proportionnelles (PSS) décrit la place du lot dans la mémoire hôte.
\chart{comparison/nvpd_focus_memory.png}{À gauche, les tampons GPU sont mesurés sur cinq paires distinctes par longueur à 2\,\%. À droite, la PSS hôte du lot est relevée après échauffement. Les durées sous les libellés donnent la moyenne et l’écart-type sur cinq blocs. Les traits fins indiquent les intervalles à 95\,\%.}{fig:nvpd-focus-memoire}{bench/repeated/divergence/focus_analysis.py}

À un million de bases et 2\,\%, le pic RSS hôte atteint en moyenne 159,5\,Mio avec NvPD en bande, 10,1\,Mio avec WFA2 et 9,6\,Mio avec Edlib. Le lot à processus unique prend 0,365\,s, contre 0,389\,s avec cinq processus.

\synthese{À un million de bases et 2\,\%, la bande ramène les tampons GPU de 22,89 à 1,59\,Mio par paire. Pour le lot de cinq paires de 100\,000 bases, un seul processus ramène la PSS hôte de 726,2 à 154,9\,Mio.}

\subsubsection{Comparaison des décompositions OpenMP à grande largeur}
Les anti-diagonales d’une matrice de \(m\times n\) cellules comptent \(m+n-1\) niveaux intérieurs de largeur variable. NvPD traite \(m\) lignes de \(n\) colonnes, dont les calculs peuvent avancer en parallèle. Les deux méthodes visitent les mêmes cellules, avec des accès mémoire et un nombre de synchronisations différents.

Le panneau a) de la figure~\ref{fig:million} compare les deux implémentations initiales sur les quinze paires d’un million de bases. Chaque méthode utilise huit cœurs EPYC. Pour les cinq paires à 5\,\% d’éditions introduites, les anti-diagonales prennent en moyenne \(190{,}24\pm0{,}42\)\,s et NvPD \(310{,}26\pm1{,}13\)\,s.

Le panneau b) reprend les mêmes cinq paires avec le noyau NvPD qui sélectionne seize prédécesseurs par permutation AVX-512. Un processus par paire dispose de huit cœurs physiques. Les cinq processus s’exécutent simultanément sur un nœud. Après une chauffe par méthode, trois blocs alternent l’ordre des deux solveurs. La lecture du FASTA précède le chronométrage, qui inclut la construction de l’index NvPD. Pour chaque paire \(j\), le gain moyen vaut
\[
  \bar S_j=\frac{1}{3}\sum_{b=1}^{3}\frac{T_{\mathrm{diag},bj}}{T_{\mathrm{NvPD},bj}} .
\]
\chart{cpu/openmp_un_million.png}{Le panneau a) donne les durées des implémentations initiales. Dans le panneau b), chaque point résume trois gains mesurés sur une paire. Le trait épais indique l’écart-type et le trait fin les extrêmes. Le losange donne la moyenne sur cinq paires, avec son écart-type et son IC à 95\,\%.}{fig:million}{bench/cpu/openmp/million/paired/trial.cpp}
Le temps moyen sur les cinq paires vaut \(196{,}864\pm0{,}123\)\,s pour les anti-diagonales et \(14{,}790\pm0{,}471\)\,s pour NvPD vectorisé. Les écarts-types décrivent la dispersion entre les cinq moyennes par paire. La figure~\ref{fig:maqao-permutation} montre comment la permutation réunit les lectures des prédécesseurs.

\synthese{À un million de bases et 5\,\% d’éditions introduites, NvPD vectorisé divise le temps des anti-diagonales par 13,322 sur huit cœurs EPYC. L’IC à 95\,\% calculé entre les cinq paires va de 12,787 à 13,857.}

